\FloatBarrier
\section{Физические границы и практический смысл результата}
Установленное преимущество бокового продольного положения относится
к заданным энергии перехода, спектру импульса и соглашению о диполе КТ.
Оно объясняется тремя связанными обстоятельствами: продольный
плазмонный максимум смещён ниже несущей; в центре боковой КТ
падающее и рассеянное поля складываются благоприятнее; реакционное
поле меньше подавляет линейный отклик, чем у торца.
Такой механизм даёт проверяемый критерий выбора геометрии через
$B$ и $K$, тогда как максимальное поле непосредственно на поверхности
не учитывает положения перехода в пространстве.

Пространственная сходимость не устраняет ограничения локальной
квазистатики. При верхней энергии основной полосы 2,234 эВ
$kc=0{,}212$; максимальное $kR$ на сетке составляет около
0,450 у торца и 0,314 сбоку. Эти параметры удовлетворяют
принятым внутренним порогам расчёта, но не являются оценкой полной
ошибки запаздывания. Отдельный диагностический расчёт длинноволновой
поправки типа Мейера--Вокауна~\cite{meier1983} оценивает сдвиг
продольной плазмонной полосы примерно в $-7$ мэВ и изменение
$|\alpha_\parallel|^2$ на несущей примерно на $-4{,}7$\%.
Эта поправка не была включена в опубликованные здесь пороги.

Для поперечного размера 9 нм существенным вопросом остаётся
поверхностное рассеяние электронов~\cite{novo2006}.
При диагностической формуле
$\Delta(\hbar\gamma)=A_s\hbar v_F/L_{\rm eff}$,
$L_{\rm eff}=4V/S=7{,}277$ нм, $v_F=1{,}4\cdot10^6$ м/с
и $A_s=0{,}3$--1 дополнительное затухание составляет
38--127 мэВ. В принятой оценке это уменьшает
$|\alpha_\parallel|^2$ на несущей примерно на 3,1--12,4\%.
Числа служат индикатором чувствительности к неучтённой физике;
они не являются рассчитанной поправкой к $\Fth$ и не задают
доверительный интервал для порога.

При малом зазоре важна и конечная протяжённость КТ.
Торцевой радиус кривизны $a^2/c=1{,}62$ нм сравним с её радиусом
2 нм. Простой показатель неоднородности дипольного поля
$3r_q/R$ при $g=0{,}5$ нм равен 0,40 у торца и 0,857 сбоку;
он не мал. Следовательно, распределённая переходная плотность
и усреднение поля по КТ могут менять количественный результат.
При $g=0{,}3$--0,5 нм также отсутствуют в модели нелокальный
отклик металла, туннелирование и перенос заряда.
Полученная численная сходимость на этих зазорах удостоверяет решение
принятых уравнений, но не оправдывает экстраполяцию к контакту.

Наконец, заданные $\gamma_1$ и $\gamma_\varphi$ не зависят от
положения относительно металла. Если после импульса $P=Q=0$,
само по себе ненулевое $\rho_{ee}$ не создаёт когерентного диполя
и не возбуждает металл через уравнение~\eqref{eq:ABK}.
Поэтому модель не содержит полного спонтанного переноса энергии,
изменения квантового выхода или геометрически зависимого
парселловского распада. Низкий порог создания заселённости не
означает автоматически усиленную фотолюминесценцию.
Для сопоставления с экспериментальным излучением потребуются
радиационные и безызлучательные каналы, параметры конкретной КТ
и реальная геометрия частицы.

\section{Заключение}
Для одиночного золотого сфероида с полуосями 12,5 и 4,5 нм
при переходе КТ 2,042 эВ наиболее эффективным на исследованной
сетке является боковое продольное возбуждение. При зазоре 0,5 нм
оно снижает порог заселённости 0,5 до $16{,}08\ \uJ$;
выигрыш относительно изолированной КТ равен 3,84, относительно
продольного торцевого положения --- 3,06. Преимущество объясняется
комплексным суммарным полем в центре КТ и обратной связью,
а не максимумом поля на поверхности. Порядок двух ведущих боковых
каналов сохраняется при всех 12 предусмотренных однофакторных
вариациях. При консервативной оценке разделимости близкую
эффективность имеют зазоры 0,5--1 нм на исходной сетке.

DD существенно искажает ближний отклик: для ведущего канала
оно занижает порог в 5,07 раза и меняет знак спектрального сдвига.
Использованная одноосцилляторная аппроксимация материала,
в свою очередь, меняет порядок двух ведущих каналов.
Согласованное многоосцилляторное материальное описание и пространственная
сходимость необходимы для количественных выводов данного сценария.

Фиксированное и оптимизированное спектральные усиления,
послеимпульсная заселённость, её порог и спектр работы поля дают
разные характеристики одной системы. Их раздельный анализ
показывает, почему оптимум монохроматического отклика может
лежать при большем зазоре, а сильное изменение узкой экситонной
структуры сопровождается лишь малым изменением полной работы.
Результаты формируют количественную основу для дальнейшего
исследования возбуждения с учётом конечного размера КТ,
поверхностного затухания и радиационных каналов.
